\documentclass[10pt,a4paper]{amsart}
\usepackage[margin=19mm]{geometry}
\usepackage{amsmath,amssymb,mathtools}
\usepackage[scr=rsfs,cal=euler]{mathalfa}
\usepackage{xfrac}
\usepackage[unicode,hidelinks,bookmarks=false]{hyperref}
\newcommand{\ie}{i.e.\ }
\newcommand{\cf}{cf.\ }
\newcommand{\resp}{respectively\ }
\newcommand{\Subsection}{\subsection}
\providecommand{\nobreakdash}{\nobreak\hbox{-}}
\setlength{\emergencystretch}{2em}
\multlinegap0pt
\usepackage{amssymb}
\usepackage{dsfont}
\usepackage[shortlabels]{enumitem}
\newtheorem{lem}{Lemma}
\title{On the structure and representations of \\ quantum graph algebras at roots of unity}
\author{St\'ephane Baseilhac, Matthieu Faitg, Philippe Roche}
\date{Source: arXiv:2601.08789v2 (CC BY 4.0)}
\begin{document}
\maketitle

\noindent\emph{Source and licence.} On the structure and representations of \\ quantum graph algebras at roots of unity. Version arXiv:2601.08789v2 (CC BY 4.0), distributed by arXiv under the Creative Commons Attribution 4.0 International licence, http://creativecommons.org/licenses/by/4.0/, as recorded in the arXiv licence metadata for this exact version and shown on the abstract page https://arxiv.org/abs/2601.08789v2. Full licence notice: \texttt{licenses/arxiv-2601.08789-CC-BY-4.0.txt}.\par\smallskip

\noindent\textit{Editorial scope.} Counted unit: the article's lem lemmeChapeau (source0.tex lines 747--758). The article's complete original proof of it is delivered with it (source0.tex lines 759--780, one \texttt{proof} environment of 22 lines). No mathematics is written, repaired or re-derived: the statement and the proof are the article's own lines, in the article's own notation, and no retained line is substituted or re-flowed.  Retained uncounted as the source-local context the proof needs: lines 438--441; lines 475--479; lines 511--513; lines 577--579; lines 626--628; lines 716--718; lines 720--722.  Nothing was omitted from any retained span, so the package carries no omission record.  No retained line contains a citation, so the wrapper carries no bibliography and no bibliography entry.  No retained line contains a figure environment, an image-inclusion command or drawing code, so no image is delivered and the package declares no asset dependency.  Editorial formatting only, in addition: an AMS \texttt{amsart} preamble that restates the article's own declarations and macros used by the retained lines, namely the environments \texttt{lem} on separate counters, together with the standard packages those lines need.  The retained environments print in this document as Lemma 1 in order of appearance, whereas the article prints the same items with the numbering of its own full document.  The complete native source of the article as distributed in the arXiv e-print for this exact version is bundled unchanged as \texttt{sources/192628/source0.tex}.
\begin{equation}\label{coregActions}
\forall \, \varphi \in H^\circ, \:\: \forall \, h,x \in H, \quad
(h \rhd \varphi)(x) = \varphi(xh), \quad (\varphi \lhd h)(x) = \varphi(hx).
\end{equation}

\begin{align}
&\sum_{(\varphi),(\psi)}\mathcal{R}(\varphi_{(1)} \otimes \psi_{(1)})\varphi_{(2)} \star \psi_{(2)} = \sum_{(\varphi),(\psi)}\mathcal{R}(\varphi_{(2)} \otimes \psi_{(2)})\psi_{(1)}\star \varphi_{(1)}, \label{braidedComm}\\
&\mathcal{R}(\varphi \star \psi \otimes \eta) = \sum_{(\eta)} \mathcal{R}\bigl(\varphi \otimes \eta_{(1)}\bigr)\mathcal{R}\bigl(\psi \otimes \eta_{(2)}\bigr),\label{coTriang1}\\
&\mathcal{R}(\varphi \otimes \psi \star \eta) = \sum_{(\varphi)} \mathcal{R}\bigl(\varphi_{(1)} \otimes \eta\bigr)\mathcal{R}\bigl(\varphi_{(2)} \otimes \psi\bigr)\label{coTriang2}
\end{align}

\begin{equation}\label{coRMatSign}
\forall \, \varphi, \psi \in H^\circ, \quad \mathcal{R}^{(\sigma)}(\varphi \otimes \psi) = \psi\bigl( \Phi^{\sigma}(\varphi) \bigr)
\end{equation}

 \begin{equation}\label{defCoactCoad}
\delta(\varphi) = \sum_{(\varphi)} \bigl( \varphi_{(1)} \star S(\varphi_{(3)}) \bigr) \otimes \varphi_{(2)}, \quad \text{written as } \delta(\varphi) = \sum_{[\varphi]} \varphi_{[1]} \otimes \varphi_{[2]}
\end{equation}

\begin{equation}\label{RSDphi}
\Phi_{0,1} : \mathcal{L}_{0,1}(H) \to H, \quad \varphi \mapsto \sum_{(\varphi)} \Phi^+(\varphi_{(1)})\Phi^-\bigl( S(\varphi_{(2)}) \bigr)
\end{equation}

\begin{equation}\label{modAlgStructHeisenberg}
(\varphi a) \cdot h = \sum_{(h)} \bigl( S(h_{(2)}) \rhd \varphi \lhd h_{(3)} \bigr) S(h_{(1)})ah_{(4)}
\end{equation}

\begin{equation}\label{HeisenbergRep}
\rho : \mathcal{H}(H^\circ) \to \mathrm{End}_\Bbbk(H^\circ), \quad \rho(\varphi a)(\psi) = \varphi \star (a \rhd \psi) \qquad (\forall \, \psi \in H^\circ).
\end{equation}

\begin{lem}\label{lemmeChapeau}
1. With the representation $\rho : \mathcal{H}(H^\circ) \to \mathrm{End}_\Bbbk(H^\circ)$ from \eqref{HeisenbergRep} we have
\begin{equation}\label{formuleRepChapeau}
\forall \, h' \in H', \:\: \forall \, \psi \in H^\circ, \quad \rho\bigl( \widehat{h'} \bigr)(\psi) = \psi \lhd S^{-1}(h').
\end{equation}
2. In $\mathcal{H}(H^\circ)$ we have the relations
\[ \textstyle \widehat{h'h''} = \widehat{h'} \, \widehat{h''}, \qquad \widehat{h'}\,\varphi a = \sum_{(h')} \bigl( \varphi \lhd S^{-1}(h'_{(2)}) \bigr)a \, \widehat{\,h'_{(1)}} \]
for all $h', h'' \in H'$, and all $\varphi a \in \mathcal{H}(H^\circ)$. In particular $\widehat{h'} \, a = a \, \widehat{h'}$ for all $a \in H \subset \mathcal{H}(H^\circ)$.
\\3. The right action \eqref{modAlgStructHeisenberg} of $h \in H$ on $\widehat{h'} \in \mathcal{H}(H^\circ)$ is given by
\[ \widehat{h'} \cdot h = \textstyle \sum_{(h)} S(h_{(1)})h_{(4)} \, \widehat{S(h_{(2)})h'h_{(3)}}. \]
4. The map $H' \to \mathcal{H}(H^\circ)$, $h' \mapsto \widehat{h'}$ is injective.
\end{lem}

\begin{proof}
1. Write $h' = \Phi_{0,1}(\varphi)$. The proof is a direct computation (using Sweedler's notation with implicit summation for coproducts and coactions):
\begin{align*}
\rho\bigl( \widehat{h'} \bigr)(\psi) &= \varphi_{[1]} \star \bigl[ S^{-1}\bigl( \Phi_{0,1}(\varphi_{[2]})  \bigr) \rhd \psi \bigr] = \bigl\langle \psi_{(2)}, S^{-1}\bigl( \Phi_{0,1}(\varphi_{(2)}) \bigr) \bigr\rangle \, \varphi_{(1)} \star S(\varphi_{(3)}) \star \psi_{(1)}\\
&=\bigl\langle S^{-1}(\psi_{(2)}), \Phi^-\bigl( S(\varphi_{(3)}) \bigr) \bigr\rangle \, \bigl\langle S^{-1}(\psi_{(3)}), \Phi^+(\varphi_{(2)}) \bigr\rangle \, \varphi_{(1)} \star S(\varphi_{(4)}) \star \psi_{(1)}\\
&=\mathcal{R}\bigl( \psi_{(2)} \otimes S(\varphi_{(3)}) \bigr) \, \mathcal{R}^{-1}\bigl( \varphi_{(2)} \otimes \psi_{(3)} \bigr) \, \varphi_{(1)} \star S(\varphi_{(4)}) \star \psi_{(1)}\\
&=\mathcal{R}\bigl( \psi_{(1)} \otimes S(\varphi_{(4)}) \bigr) \, \mathcal{R}^{-1}\bigl( \varphi_{(2)} \otimes \psi_{(3)} \bigr) \, \varphi_{(1)} \star \psi_{(2)} \star S(\varphi_{(3)})\\
&=\mathcal{R}\bigl( \psi_{(1)} \otimes S(\varphi_{(4)}) \bigr) \, \mathcal{R}^{-1}\bigl( \varphi_{(1)} \otimes \psi_{(2)} \bigr) \, \psi_{(3)} \star \varphi_{(2)} \star S(\varphi_{(3)})\\
&=\mathcal{R}\bigl( \psi_{(1)} \otimes S(\varphi_{(2)}) \bigr) \, \mathcal{R}^{-1}\bigl( \varphi_{(1)} \otimes \psi_{(2)} \bigr) \, \psi_{(3)}\\
&=\bigl\langle S^{-1}(\psi_{(1)}), \Phi^-\bigl( S(\varphi_{(2)}) \bigr) \bigr\rangle \, \bigl\langle S^{-1}(\psi_{(2)}), \Phi^+(\varphi_{(1)}) \bigr\rangle \, \psi_{(3)}\\
&=\bigl\langle \psi_{(1)}, S^{-1}\bigl[ \Phi^+(\varphi_{(1)}) \Phi^-\bigl( S(\varphi_{(2)}) \bigr) \bigr] \bigr\rangle \, \psi_{(2)} = \psi \lhd S^{-1}\bigl( \Phi_{0,1}(\varphi) \bigr) = \psi \lhd S^{-1}(h')
\end{align*}
where the first equality is by definition of $\rho$, the second equality is by definition of the coaction \eqref{defCoactCoad} and of $\rhd$ in \eqref{coregActions}, the third is by definition of $\Phi_{0,1}$ in \eqref{RSDphi}, the fourth is by \eqref{coRMatSign}, the fifth and sixth use \eqref{braidedComm}, and the eighth is by \eqref{coRMatSign} also using that $\mathcal{R}^{-1} = \mathcal{R} \circ (\mathrm{id} \otimes S^{-1})$ and the remaining equalities are easy.
\\2. These claims are easily deduced from item 1, thanks to injectivity of $\rho$. For instance:
\begin{align*}
\rho\bigl( \widehat{h'h''} \bigr)(\psi) = \psi \lhd S^{-1}(h'h'') &= \bigl( \psi \lhd S^{-1}(h'') \bigr) \lhd S^{-1}(h')\\
&= \bigl[ \rho\bigl(\widehat{h'}\bigr) \circ \rho\bigl( \widehat{h''} \bigr) \bigr](\psi) = \rho\bigl( \widehat{h'}\,\widehat{h''} \bigr)(\psi)
\end{align*}
for all $\psi \in H^\circ$; hence $\rho\bigl( \widehat{h'h''} \bigr) = \rho\bigl( \widehat{h'}\,\widehat{h''} \bigr)$ and $\widehat{h'h''} = \widehat{h'}\,\widehat{h''}$.
\\3. It suffices to check this equality through the injective morphism $\rho : \mathcal{H}(H^\circ) \to \mathrm{End}_\Bbbk(H^\circ)$, which is a straightforward computation using \eqref{formuleRepChapeau}.
\\4. Recall that $\mathcal{H}(H^\circ) = H^\circ \otimes H$ as a vector space. By definition of a coaction it holds $\varepsilon_{H^\circ}(\varphi_{[1]})\varphi_{[2]} = \varphi$. As a result writing $h' = \Phi_{0,1}(\varphi)$ we find $(\varepsilon_{H^\circ} \otimes S)\bigl(\widehat{h'} \bigr) = h'$.
\end{proof}

\end{document}
